\section{Discussion}

\paragraph{How does model behavior change with scale?}
All families improve with more labels, but TF-IDF reaches useful performance earlier. Learned-from-scratch neural models close the gap near the upper end of each schedule. NSMC is the only domain where the neural advantage reaches and sustains the prespecified one-point margin. Shopping's best neural pipeline wins at full data without clearing the margin; Steam's top three neural methods are tightly grouped and also remain below it.

\paragraph{Does the improvement justify cost?}
There is no universal answer. Linear TF-IDF is fast, compact, and competitive at every tested scale. On NSMC, a team that values roughly \Result{story_full_gap_nsmc} additional macro-F1 points may accept \Result{story_fit_ratio_nsmc} times the fit time. On Shopping and Steam, the smaller gains make that trade harder to justify unless predictive quality dominates deployment constraints. Random Forest is rarely attractive because its sparse-text latency and model size are high without a matching quality gain.

\paragraph{What transfers across domains?}
The shared result is that pooling existing source labels substantially improves a low-label target-only Transformer. The domain-specific result is whether that pooled model beats TF-IDF and whether a contrastive loss helps. Shopping remains especially favorable to TF-IDF at low target budgets. Standard SupCon and DB-CDSupCon each win selected cells, but neither dominates pooled CE across targets and sizes.

\paragraph{Does DB-CDSupCon improve representations?}
The separation diagnostic changes in the intended direction only on Shopping when averaged across target sizes. NSMC and Steam move slightly downward. Even where separation rises, target macro-F1 does not rise at every size. A representation diagnostic is supporting evidence, not a substitute for the task metric.

\paragraph{Practical recommendations.}
With hundreds or a few thousand labels, start with TF-IDF plus a linear classifier. At tens of thousands of labels, compare that baseline with a small Transformer or gated RNN using the target hardware. If labeled source domains already exist, pooled CE is the essential transfer baseline. Add a contrastive objective only after a matched pooled comparison shows a stable benefit for the target domain.
